\documentclass[12pt,a4paper]{article}

\usepackage[T2A]{fontenc}
\usepackage[utf8]{inputenc}
\usepackage[russian]{babel}
\usepackage{geometry}
\geometry{a4paper, margin=2.5cm}
\usepackage{hyperref}
\hypersetup{
    colorlinks=true,
    linkcolor=blue,
    urlcolor=blue,
    citecolor=blue,
}
\usepackage{listings}
\usepackage{xcolor}
\usepackage{graphicx}
\usepackage{titlesec}
\usepackage{setspace}

\lstset{
    basicstyle=\ttfamily\small,
    breaklines=true,
    breakatwhitespace=false,
    columns=fullflexible,
    frame=single,
    framesep=4pt,
    framerule=0.4pt,
    rulecolor=\color{gray!50},
    backgroundcolor=\color{gray!5},
    showstringspaces=false,
    keepspaces=true,
    extendedchars=true,
    inputencoding=utf8,
    literate=
        {а}{{\selectfont\char224}}1 {б}{{\selectfont\char225}}1
        {в}{{\selectfont\char226}}1 {г}{{\selectfont\char227}}1
        {д}{{\selectfont\char228}}1 {е}{{\selectfont\char229}}1
        {ё}{{\selectfont\char184}}1 {ж}{{\selectfont\char230}}1
        {з}{{\selectfont\char231}}1 {и}{{\selectfont\char232}}1
        {й}{{\selectfont\char233}}1 {к}{{\selectfont\char234}}1
        {л}{{\selectfont\char235}}1 {м}{{\selectfont\char236}}1
        {н}{{\selectfont\char237}}1 {о}{{\selectfont\char238}}1
        {п}{{\selectfont\char239}}1 {р}{{\selectfont\char240}}1
        {с}{{\selectfont\char241}}1 {т}{{\selectfont\char242}}1
        {у}{{\selectfont\char243}}1 {ф}{{\selectfont\char244}}1
        {х}{{\selectfont\char245}}1 {ц}{{\selectfont\char246}}1
        {ч}{{\selectfont\char247}}1 {ш}{{\selectfont\char248}}1
        {щ}{{\selectfont\char249}}1 {ъ}{{\selectfont\char250}}1
        {ы}{{\selectfont\char251}}1 {ь}{{\selectfont\char252}}1
        {э}{{\selectfont\char253}}1 {ю}{{\selectfont\char254}}1
        {я}{{\selectfont\char255}}1
        {А}{{\selectfont\char192}}1 {Б}{{\selectfont\char193}}1
        {В}{{\selectfont\char194}}1 {Г}{{\selectfont\char195}}1
        {Д}{{\selectfont\char196}}1 {Е}{{\selectfont\char197}}1
        {Ё}{{\selectfont\char168}}1 {Ж}{{\selectfont\char198}}1
        {З}{{\selectfont\char199}}1 {И}{{\selectfont\char200}}1
        {Й}{{\selectfont\char201}}1 {К}{{\selectfont\char202}}1
        {Л}{{\selectfont\char203}}1 {М}{{\selectfont\char204}}1
        {Н}{{\selectfont\char205}}1 {О}{{\selectfont\char206}}1
        {П}{{\selectfont\char207}}1 {Р}{{\selectfont\char208}}1
        {С}{{\selectfont\char209}}1 {Т}{{\selectfont\char210}}1
        {У}{{\selectfont\char211}}1 {Ф}{{\selectfont\char212}}1
        {Х}{{\selectfont\char213}}1 {Ц}{{\selectfont\char214}}1
        {Ч}{{\selectfont\char215}}1 {Ш}{{\selectfont\char216}}1
        {Щ}{{\selectfont\char217}}1 {Ъ}{{\selectfont\char218}}1
        {Ы}{{\selectfont\char219}}1 {Ь}{{\selectfont\char220}}1
        {Э}{{\selectfont\char221}}1 {Ю}{{\selectfont\char222}}1
        {Я}{{\selectfont\char223}}1
}

\begin{document}

% ============================================================
% Титульный лист
% ============================================================
\begin{titlepage}
    \centering
    {\normalsize Санкт-Петербургский политехнический университет Петра Великого\par}
    {\normalsize Институт компьютерных наук и кибербезопасности\par}
    {\normalsize Высшая школа программной инженерии\par}

    \vspace{4cm}

    {\Huge\bfseries Практическое задание №5\par}

    \vspace{0.5cm}
    {\large по дисциплине <<Системы программирования>>\par}

    \vfill

    \begin{flushleft}
        \noindent Выполнил:\hfill Подшивалов Г. А.\\[0.3em]
        \noindent Группа:\hfill гр. 5130904/50501\\[0.3em]
        \noindent Проверила:\hfill Степина Н. О.
    \end{flushleft}

    \vfill

    {\normalsize Санкт-Петербург\par}
    {\normalsize 2026\par}
\end{titlepage}

% ============================================================
% Содержание
% ============================================================
\tableofcontents
\newpage

% ============================================================
% 1. Постановка задачи
% ============================================================
\section{Постановка задачи}

Цель данного задания — продемонстрировать использование Ansible для решения задачи удалённого развертывания среды эксплуатации ИТ-системы. В качестве облачной платформы выбран Yandex Cloud. Необходимо реализовать четыре сценария интеграции:

\begin{itemize}
    \item provisioning виртуальных машин через плейбук Ansible;
    \item динамический инвентарь (Ansible автоматически находит созданные ВМ в облаке);
    \item configuration management (установка и настройка nginx);
    \item teardown — корректное удаление созданных ресурсов.
\end{itemize}

% ============================================================
% 2. Ход работы
% ============================================================
\section{Ход работы}

\subsection{Подготовка инфраструктуры}

\subsubsection{Установка Yandex Cloud CLI}

Выполнена установка Yandex Cloud CLI на локальной Windows-машине:

\begin{lstlisting}
iex (New-Object System.Net.WebClient).DownloadString(
    'https://storage.yandexcloud.net/yandexcloud-yc/install.ps1')
\end{lstlisting}

\begin{lstlisting}
yc --version
Yandex Cloud CLI 0.150.0 windows/amd64
\end{lstlisting}

\subsubsection{Инициализация CLI и создание сервисного аккаунта}

Выполнена настройка профиля, создание каталога и сервисного аккаунта для Ansible:

\begin{lstlisting}
yc init
yc iam service-account create --name ansible-sa
yc resource-manager folder add-access-binding lab5-folder `
    --role editor `
    --subject serviceAccount:<SA_ID>
yc iam key create --service-account-name ansible-sa `
    --output $HOME\.yc\key.json
\end{lstlisting}

\subsubsection{Создание SSH-ключей}

\begin{lstlisting}
ssh-keygen -t ed25519 -f $HOME\.ssh\lab5_key -N '""'
\end{lstlisting}

\subsubsection{Создание управляющей ВМ (control-node)}

\begin{lstlisting}
yc compute instance create `
    --name control-node `
    --zone ru-central1-a `
    --platform-id standard-v3 `
    --cores 2 --memory 2 `
    --create-boot-disk image-folder-id=standard-images,image-family=ubuntu-2204-lts,size=20 `
    --network-interface subnet-name=default-ru-central1-a,nat-ip-version=ipv4 `
    --ssh-key $HOME\.ssh\lab5_key.pub `
    --labels role=control,lab=lab5
\end{lstlisting}

\begin{lstlisting}
id: fhmc34s7ns0h6stgnb6m
name: control-node
status: RUNNING
...
primary_v4_address:
    address: 192.168.1.13
    one_to_one_nat:
        address: 84.201.132.241
\end{lstlisting}

\subsubsection{Подключение к ВМ}

\begin{lstlisting}
ssh -i ~/.ssh/lab5_key yc-user@84.201.132.241
\end{lstlisting}

\begin{lstlisting}
Welcome to Ubuntu 22.04.5 LTS (GNU/Linux 5.15.0-176-generic x86_64)
yc-user@control-node:~$
\end{lstlisting}

\subsection{Установка Ansible и зависимостей}

\subsubsection{Обновление системы и установка пакетов}

\begin{lstlisting}
sudo apt update && sudo apt upgrade -y
sudo apt install -y python3 python3-pip git
\end{lstlisting}

\subsubsection{Установка Python-библиотек}

\begin{lstlisting}
pip3 install --user ansible yandexcloud
echo 'export PATH=$HOME/.local/bin:$PATH' >> ~/.bashrc
source ~/.bashrc
\end{lstlisting}

\begin{lstlisting}
ansible --version
ansible [core 2.17.0]
\end{lstlisting}

\subsubsection{Установка yc CLI на control-node}

\begin{lstlisting}
curl -sSL https://storage.yandexcloud.net/yandexcloud-yc/install.sh | bash
source ~/.bashrc
yc config profile create sa-profile
yc config set service-account-key ~/.yc/key.json
yc config set cloud-id <CLOUD_ID>
yc config set folder-id <FOLDER_ID>
yc config set compute-default-zone ru-central1-a
\end{lstlisting}

\subsection{Динамический инвентарь}

\subsubsection{Создание скрипта инвентаря}

Создан скрипт \texttt{inventory/yc\_fallback.py}, который вызывает \texttt{yc compute instance list --format json}, фильтрует запущенные ВМ и группирует их по меткам (\texttt{labels}). Группировка: ВМ с меткой \texttt{role=web} попадают в группу \texttt{role\_web}.

\begin{lstlisting}
chmod +x ~/lab5/inventory/yc_fallback.py
~/lab5/inventory/yc_fallback.py | python3 -m json.tool
\end{lstlisting}

\begin{lstlisting}
{
    "_meta": {
        "hostvars": {
            "control-node": {
                "ansible_host": "84.201.132.241",
                "ansible_user": "yc-user"
            }
        }
    },
    "role_control": {"hosts": ["control-node"]},
    "lab_lab5":    {"hosts": ["control-node"]}
}
\end{lstlisting}

\subsubsection{Конфигурация Ansible}

Скрипт прописан как инвентарь по умолчанию в \texttt{ansible.cfg}:

\begin{lstlisting}
[defaults]
inventory            = ./inventory/yc_fallback.py
roles_path           = ./roles
host_key_checking    = False
remote_user          = yc-user
private_key_file     = ~/.ssh/lab5_key
stdout_callback      = yaml
\end{lstlisting}

\subsubsection{Тестирование динамического инвентаря}

\begin{lstlisting}
ansible-inventory --graph
@all:
  |--@ungrouped:
  |--@role_control:
  |  |--control-node
  |--@lab_lab5:
  |  |--control-node
\end{lstlisting}

\subsection{Создание ВМ через Ansible Playbook}

\subsubsection{Playbook provision.yml}

Создан плейбук, который через \texttt{yc} CLI разворачивает две ВМ \texttt{web-1} и \texttt{web-2} с метками \texttt{role=web, lab=lab5}, дожидается готовности SSH и регистрирует их в in-memory инвентаре.

Ключевые задачи:

\begin{lstlisting}
- name: "Создать ВМ ({{ item.name }})"
  ansible.builtin.command:
    cmd: >
      yc compute instance create
      --name {{ item.name }}
      --zone {{ yc.zone }}
      --platform-id {{ vm_defaults.platform_id }}
      --cores {{ vm_defaults.cores }}
      --memory {{ vm_defaults.memory_gb }}
      --create-boot-disk image-folder-id={{ image.folder_id }},image-family={{ image.family }},size={{ vm_defaults.disk_size_gb }}
      --network-interface subnet-id={{ yc.subnet_id }},nat-ip-version=ipv4
      --ssh-key {{ lookup('env', 'HOME') }}/.ssh/lab5_key.pub
      --labels role=web,lab=lab5
      --format json
  loop: "{{ web_hosts }}"
  when: item.name not in existing_names

- name: "Подождать готовности SSH"
  ansible.builtin.wait_for:
    host: "{{ item.network_interfaces[0].primary_v4_address.one_to_one_nat.address }}"
    port: 22
    delay: 5
    timeout: 300
  loop: "{{ web_inventory }}"
\end{lstlisting}

\subsubsection{Запуск playbook}

\begin{lstlisting}
ansible-playbook playbooks/provision.yml
\end{lstlisting}

\begin{lstlisting}
PLAY [Provision web VMs in Yandex Cloud] ***********************

TASK [Получить список уже существующих ВМ] *********************
ok: [localhost]

TASK [Создать ВМ (web-1)] **************************************
changed: [localhost] => (item=web-1)

TASK [Создать ВМ (web-2)] **************************************
changed: [localhost] => (item=web-2)

TASK [Подождать готовности SSH на каждой web-ВМ] ***************
ok: [localhost] => (item=web-1)
ok: [localhost] => (item=web-2)

TASK [Вывести результат] ***************************************
ok: [localhost] => msg: "VM web-1 -> 89.169.151.155"
ok: [localhost] => msg: "VM web-2 -> 93.77.191.135"

PLAY RECAP ******************************************************
localhost : ok=7 changed=2 unreachable=0 failed=0
\end{lstlisting}

\subsection{Настройка nginx на ВМ}

\subsubsection{Роль webserver}

Создана Ansible-роль \texttt{webserver}, устанавливающая nginx и публикующая шаблонную страницу со сведениями о хосте (\texttt{ansible\_hostname}, IP, ОС, ядро, CPU, RAM).

\begin{lstlisting}
# roles/webserver/tasks/main.yml
- name: "Установить nginx"
  ansible.builtin.apt:
    name: nginx
    state: present
    update_cache: true

- name: "Развернуть стартовую страницу"
  ansible.builtin.template:
    src: index.html.j2
    dest: /var/www/html/index.nginx-debian.html
    owner: root
    group: root
    mode: "0644"
  notify: "Restart nginx"

- name: "Включить и запустить nginx"
  ansible.builtin.service:
    name: nginx
    state: started
    enabled: true
\end{lstlisting}

\subsubsection{Playbook configure.yml}

\begin{lstlisting}
- name: "Configure web servers (nginx)"
  hosts: role_web
  become: true
  pre_tasks:
    - name: "Подождать SSH"
      ansible.builtin.wait_for_connection:
        timeout: 120
  roles:
    - role: webserver
\end{lstlisting}

\subsubsection{Запуск playbook}

\begin{lstlisting}
ansible-playbook playbooks/configure.yml
\end{lstlisting}

\begin{lstlisting}
PLAY [Configure web servers (nginx)] ***************************

TASK [Gathering Facts] *****************************************
ok: [web-1]
ok: [web-2]

TASK [webserver : Установить nginx] ****************************
changed: [web-1]
changed: [web-2]

TASK [webserver : Развернуть стартовую страницу] ***************
changed: [web-1]
changed: [web-2]

TASK [webserver : Включить и запустить nginx] ******************
ok: [web-1]
ok: [web-2]

RUNNING HANDLER [webserver : Restart nginx] ********************
changed: [web-1]
changed: [web-2]

PLAY RECAP *****************************************************
web-1 : ok=6 changed=3 unreachable=0 failed=0
web-2 : ok=6 changed=3 unreachable=0 failed=0
\end{lstlisting}

\subsubsection{Проверка работы веб-сервера}

\begin{lstlisting}
curl http://89.169.151.155
\end{lstlisting}

\begin{lstlisting}
<!doctype html>
<html lang="ru">
<head>
    <meta charset="utf-8">
    <title>Lab 5 — web-1</title>
</head>
<body>
    <h1>Лабораторная работа №5</h1>
    <p>Эта страница развёрнута Ansible на ВМ <code>web-1</code>
       в Yandex Cloud.</p>
    <table>
        <tr><th>Hostname</th><td>web-1</td></tr>
        <tr><th>Внутренний IPv4</th><td>192.168.1.21</td></tr>
        <tr><th>ОС</th><td>Ubuntu 22.04</td></tr>
        <tr><th>Ядро</th><td>5.15.0-176-generic</td></tr>
        <tr><th>CPU</th><td>2 vCPU</td></tr>
        <tr><th>RAM, МБ</th><td>1990</td></tr>
    </table>
</body>
</html>
\end{lstlisting}

\subsubsection{Проверка идемпотентности}

Повторный запуск \texttt{configure.yml} не должен вносить изменений:

\begin{lstlisting}
ansible-playbook playbooks/configure.yml

PLAY RECAP *****************************************************
web-1 : ok=6 changed=0 unreachable=0 failed=0
web-2 : ok=6 changed=0 unreachable=0 failed=0
\end{lstlisting}

\subsection{Удаление ВМ через Ansible}

\subsubsection{Playbook teardown.yml}

Плейбук удаляет ВМ с метками \texttt{lab=lab5} и \texttt{role=web}, control-node при этом не затрагивается:

\begin{lstlisting}
- name: "Teardown web VMs in Yandex Cloud"
  hosts: localhost
  tasks:
    - name: "Получить актуальный список ВМ"
      ansible.builtin.command:
        cmd: yc compute instance list --format json
      register: list_raw

    - name: "Выбрать имена web-ВМ для удаления"
      ansible.builtin.set_fact:
        targets: >-
          {{ (list_raw.stdout | from_json)
             | selectattr('labels.lab', 'equalto', 'lab5')
             | selectattr('labels.role', 'equalto', 'web')
             | map(attribute='name') | list }}

    - name: "Удалить ВМ"
      ansible.builtin.command:
        cmd: "yc compute instance delete --name {{ item }}"
      loop: "{{ targets }}"
\end{lstlisting}

\subsubsection{Запуск playbook}

\begin{lstlisting}
ansible-playbook playbooks/teardown.yml
\end{lstlisting}

\begin{lstlisting}
PLAY [Teardown web VMs in Yandex Cloud] ************************

TASK [Получить актуальный список ВМ] ***************************
ok: [localhost]

TASK [Показать что будет удалено] ******************************
ok: [localhost] => msg: "Будут удалены ВМ: web-1, web-2"

TASK [Удалить ВМ] **********************************************
changed: [localhost] => (item=web-1)
changed: [localhost] => (item=web-2)

PLAY RECAP *****************************************************
localhost : ok=4 changed=2 unreachable=0 failed=0
\end{lstlisting}

После выполнения в каталоге остаётся только \texttt{control-node}:

\begin{lstlisting}
yc compute instance list
+----------------------+--------------+---------------+---------+
|          ID          |     NAME     |     ZONE      | STATUS  |
+----------------------+--------------+---------------+---------+
| fhmc34s7ns0h6stgnb6m | control-node | ru-central1-a | RUNNING |
+----------------------+--------------+---------------+---------+
\end{lstlisting}

% ============================================================
% 3. Вывод
% ============================================================
\section{Вывод}

В ходе выполнения лабораторной работы были успешно освоены следующие навыки:

\begin{itemize}
    \item работа с Yandex Cloud CLI и создание облачной инфраструктуры;
    \item установка и настройка Ansible на управляющей виртуальной машине;
    \item создание динамического инвентаря для автоматического обнаружения ВМ через API облака;
    \item автоматизация развёртывания виртуальных машин с помощью Ansible playbook;
    \item установка и настройка веб-сервера nginx через роли Ansible с использованием Jinja2-шаблонов;
    \item обеспечение идемпотентности конфигурации (повторный запуск playbook не вносит изменений);
    \item удаление виртуальных машин средствами Ansible с фильтрацией по меткам.
\end{itemize}

Продемонстрирована применимость подхода Infrastructure as Code (IaC) для интеграции Ansible с облачной платформой Yandex Cloud. Все цели лабораторной работы достигнуты, инфраструктура протестирована и работает корректно.

\end{document}
